% answers_short.tex —— 2013 年·安徽卷·文科 · 简短答案速查
% 由 tools/make_short_answers.py 自 answers.tex 的 \gksolhead/\gkanswer 提取，
% 并与 2013/2013_数学真题_安徽文_答案要点.md 逐题比对后生成，勿手改（改源 MD 后重跑）。

\begin{gknotice}
\textbf{使用说明}：本节只给最终答案便于快速核对，完整推导见解析版。答案为编写者独立推导，\textbf{非官方标准答案}。
\end{gknotice}

\gkpart{选择题}
\begin{center}\begin{tabular}{ccccccccccc}
\toprule
题号 & 1 & 2 & 3 & 4 & 5 & 6 & 7 & 8 & 9 & 10 \\
\midrule
答案 & D & A & C & B & D & C & A & B & B & A \\
\bottomrule
\end{tabular}\end{center}

\gkpart{填空题}
\begin{itemize}[leftmargin=2.6em]
  \item[\textbf{11.}] $(0,1]$。
  \item[\textbf{12.}] $\dfrac73$。
  \item[\textbf{13.}] $-\dfrac13$。
  \item[\textbf{14.}] $-\dfrac{x(x+1)}2$。
  \item[\textbf{15.}] ①②③⑤。
\end{itemize}

\gkpart{解答题}
\begin{itemize}[leftmargin=2.6em]
  \item[\textbf{16.}] (1) 最小值为 $-\sqrt3$，此时 $x$ 的集合为 $\left\{x\ \middle|\ x=2k\pi-\dfrac{2\pi}3,\ k\in\mathbf Z\right\}$；(2) 将 $y=\sin x$ 的图象向左平移 $\dfrac\pi6$ 个单位，再把所得图象上各点的纵坐标伸长为原来的 $\sqrt3$ 倍（横坐标不变），即得 $y=f(x)$ 的图象
  \item[\textbf{17.}] (1) 甲校高三年级学生总人数为 $600$，及格率估计为 $\dfrac56$；(2) $\overline{x}_1-\overline{x}_2=0.5$
  \item[\textbf{18.}] (1) 证明见解析；(2) $\dfrac{\sqrt2}6$
  \item[\textbf{19.}] (1) $a_n=n+1$（$n\in\mathbf N^*$）；(2) $S_n=n^2+3n+1-\dfrac1{2^n}$
  \item[\textbf{20.}] (1) $I$ 的长度为 $\dfrac a{1+a^2}$；(2) 最小值为 $\dfrac{1-k}{k^2-2k+2}$
  \item[\textbf{21.}] (1) $\dfrac{x^2}8+\dfrac{y^2}4=1$；(2) 一定有唯一的公共点，$QG$ 是椭圆在 $Q$ 处的切线，证明见解析
\end{itemize}
